\documentclass[11pt,a4paper]{article}

% ---------------------------------------------------------
% PREAMBLE (stable, warning-free, Overleaf-safe)
% ---------------------------------------------------------
\usepackage[utf8]{inputenc}
\usepackage[T1]{fontenc}
\usepackage{lmodern}
\usepackage{amsmath, amssymb, amsfonts}
\usepackage{bm}
\usepackage{geometry}
\usepackage{graphicx}
\usepackage{hyperref}
\usepackage{cite}
\usepackage{authblk}
\usepackage{physics}
\usepackage{mathtools}
\usepackage{enumitem}

\geometry{margin=2.4cm}

\title{\textbf{Companion XXXII: Halo Concentration, Density Profiles, and Substructure in the Relaxation-Driven Cyclic Cosmology}}
\author{Michael Lehmann \\ \texttt{mi.lehmann@gmx.de}}
\date{April 2026}

\begin{document}
\maketitle

\begin{abstract}
This Companion investigates how the Relaxation-Driven Cyclic Cosmology (RDCC) modifies 
the internal structure of dark matter halos. While Companion~XXXI established the 
non-linear halo model and its global implications, the present work focuses on the 
microstructure of halos: their density profiles, concentration--mass relation, 
subhalo abundance, and the emergence of core-like features. The relaxation-induced 
suppression of small-scale potential evolution alters the collapse history of halos 
and reshapes their internal architecture. We derive analytical estimates, discuss 
scaling relations, and outline observational signatures relevant for weak lensing, 
galaxy--galaxy lensing, and cluster astrophysics.
\end{abstract}
% ---------------------------------------------------------
\section{Introduction}

Dark matter halos are not only the building blocks of cosmic structure but also 
repositories of dynamical information. Their internal density profiles, concentration 
parameters, and substructure content encode the history of gravitational collapse. 
In the standard cosmological model, halo structure is shaped by hierarchical growth 
and the scale-free nature of cold dark matter.

In the Relaxation-Driven Cyclic Cosmology (RDCC), the situation changes 
fundamentally. The relaxation-modified potential evolution introduces a 
scale-dependent suppression of small-scale growth, which alters the timing and 
efficiency of halo collapse. As a result, halo profiles, concentrations, and 
subhalo populations acquire characteristic RDCC signatures.

This Companion develops the theoretical framework for halo microstructure in RDCC.
% ---------------------------------------------------------
\section{Density Profiles in RDCC}

The standard Navarro–Frenk–White (NFW) profile is
\begin{equation}
    \rho_{\mathrm{NFW}}(r)
    = \frac{\rho_{s}}{(r/r_{s})(1+r/r_{s})^{2}}.
\end{equation}

In RDCC, the suppression of small-scale potential evolution modifies the inner 
collapse dynamics. We model the RDCC profile as
\begin{equation}
    \rho_{\mathrm{RDCC}}(r)
    = \rho_{\mathrm{NFW}}(r)
      \left[1 - \chi_{\rho}
      \left(\frac{r}{r_{\mathrm{rel}}}\right)^{\beta}\right],
\end{equation}
where $r_{\mathrm{rel}}$ is the scale corresponding to $k_{\mathrm{rel}}$.

RDCC predicts:
\begin{itemize}[leftmargin=1.2cm]
    \item slightly shallower inner slopes,
    \item smoother transitions between inner and outer regions,
    \item reduced central densities for low-mass halos,
    \item preserved outer profiles for massive halos.
\end{itemize}
% ---------------------------------------------------------
\section{Concentration--Mass Relation}

The concentration parameter is defined as
\begin{equation}
    c(M) = \frac{r_{\mathrm{vir}}}{r_{s}}.
\end{equation}

In RDCC, the delayed collapse of small halos and the earlier collapse of massive 
halos modify the concentration--mass relation:
\begin{equation}
    c_{\mathrm{RDCC}}(M)
    = c(M)
      \left[1 - \chi_{c}
      \left(\frac{k(M)}{k_{\mathrm{rel}}}\right)^{\alpha}\right].
\end{equation}

Consequences:
\begin{itemize}[leftmargin=1.2cm]
    \item reduced concentrations for dwarf-scale halos,
    \item nearly unchanged concentrations for cluster-scale halos,
    \item a characteristic turnover mass $M_{\mathrm{rel}}$,
    \item smoother redshift evolution of $c(M)$.
\end{itemize}
% ---------------------------------------------------------
\section{Subhalo Abundance and Survival}

Subhalos trace the hierarchical assembly of halos. Their abundance is commonly 
described by the subhalo mass function:
\begin{equation}
    \frac{dN}{d\ln m}
    \propto m^{-\gamma}.
\end{equation}

In RDCC, the suppression of small-scale growth reduces the formation of low-mass 
subhalos:
\begin{equation}
    \left(\frac{dN}{d\ln m}\right)_{\mathrm{RDCC}}
    = \left(\frac{dN}{d\ln m}\right)
      \left[1 - \chi_{\mathrm{sub}}
      \left(\frac{k(m)}{k_{\mathrm{rel}}}\right)^{\alpha}\right].
\end{equation}

RDCC predicts:
\begin{itemize}[leftmargin=1.2cm]
    \item fewer low-mass subhalos,
    \item enhanced survival of intermediate-mass subhalos,
    \item reduced tidal stripping due to smoother potentials,
    \item a natural explanation for missing-satellite-like features.
\end{itemize}
% ---------------------------------------------------------
\section{Core Formation and Inner Structure}

The relaxation-modified potential evolution leads to a reduced efficiency of 
small-scale collapse. This can produce core-like features in low-mass halos.

We model the RDCC core radius as
\begin{equation}
    r_{\mathrm{core}}
    \sim r_{s}
    \left(\frac{k_{\mathrm{rel}}}{k(M)}\right)^{\alpha}.
\end{equation}

Implications:
\begin{itemize}[leftmargin=1.2cm]
    \item dwarf-scale halos develop shallow cores,
    \item Milky-Way-scale halos remain cusp-like,
    \item cluster-scale halos are unaffected,
    \item core sizes correlate with $k_{\mathrm{rel}}$.
\end{itemize}
% ---------------------------------------------------------
\section{Observational Signatures}

The RDCC modifications to halo structure affect several observables:

\begin{itemize}[leftmargin=1.2cm]
    \item \textbf{Galaxy--galaxy lensing:} reduced small-scale shear signals.
    \item \textbf{Weak lensing peaks:} fewer high-concentration peaks.
    \item \textbf{Rotation curves:} natural emergence of shallow cores.
    \item \textbf{Cluster lensing:} preserved NFW-like outer profiles.
    \item \textbf{Substructure lensing:} reduced subhalo-induced perturbations.
\end{itemize}

These signatures provide a direct test of RDCC on non-linear scales.
% ---------------------------------------------------------
\section{Conclusion}

This Companion develops the internal structure of dark matter halos in the 
Relaxation-Driven Cyclic Cosmology. The relaxation-modified potential evolution 
reshapes halo profiles, concentrations, and subhalo populations. These effects 
propagate into lensing, galaxy dynamics, and the morphology of the cosmic web. 
Together with Companion~XXXI, the present work establishes the foundation for 
RDCC non-linear cosmology and prepares the ground for future studies of galaxy 
formation, cluster astrophysics, and cosmic web topology.

% ========================
% References
% ========================

\begin{thebibliography}{10}

\bibitem{Flagship} Lehmann, M. (2026). Relaxation-Driven Cyclic Cosmology (RDCC): A Minimal CPT-Symmetric Framework with a Single Parameter. Flagship Paper. Zenodo: \url{https://zenodo.org/records/19744958} (v43.0 or later).

\bibitem{Zenodo} The complete RDCC companion series (I--LVII) is available at the same Zenodo record above.

% Thematisch relevante Companions
\bibitem{CompXVI} Companion XVI: Boltzmann Code and Modified Hierarchy.
\bibitem{CompXXXI} Companion XXXI: RDCC Halo Model and Non-Linear Structure.
\bibitem{CompXXXII} Companion XXXII--XXXIX: Galaxy--Halo Connection, Cosmic Web, Lensing, RSD, ISW, tSZ, Rotation Curves, CGM.

% Wichtige externe Referenzen
\bibitem{Planck2020} Planck Collaboration (2020), Astron. Astrophys. 641, A6.
\bibitem{DESI2024} DESI Collaboration (2024/2025), arXiv: [aktuelle $f\sigma_8$/BAO-Paper].
\bibitem{JWST2024} Maiolino et al. (2024), Nature (JWST high-$z$ galaxies/quasars).
\bibitem{Kaiser1987} Kaiser, N. (1987), Mon. Not. R. Astron. Soc. 227, 1 (RSD).
\bibitem{Bartelmann2001} Bartelmann, M. \& Schneider, P. (2001), Phys. Rep. 340, 291 (Weak Lensing Review).
\bibitem{HuOkamoto2002} Hu, W. \& Okamoto, T. (2002), Astrophys. J. 574, 566 (CMB Lensing).
\bibitem{LewisChallinor2006} Lewis, A. \& Challinor, A. (2006), Phys. Rep. 429, 1 (CMB Lensing Review).
\bibitem{NFW1997} Navarro, J. F., Frenk, C. S. \& White, S. D. M. (1997), Astrophys. J. 490, 493 (Halo Profiles).

\end{thebibliography}

\end{document}
